@ID: OW23D1P1I
@BIDANG: Analisis Real

Misalkan fungsi $f,g:[a,b]\to\mathbb R$ memiliki turunan hingga orde $n-1$ yang kontinu pada $[a,b]$. Jika untuk suatu $c\in[a,b]$ berlaku

$$
f^{(k)}(c)=0,\qquad
g^{(k)}(c)=0,\qquad k=0,1,\ldots,n-2,
$$

serta

$$
g^{(n-1)}(c)\ne0,
$$

tentukan

$$
\lim_{x\to c}\frac{e^xg(x)}{f(x)}.
$$

@ID: OW23D1P2I
@BIDANG: Analisis Real

Diberikan barisan bilangan real $\{x_n\}$ dan $\{y_n\}$. Jika

$$
x_n=\sum_{k=1}^{n}\frac{(-1)^k}{(2k)!}
$$

dan

$$
y_n=1+x_n,
$$

tentukan

$$
\lim_{n\to\infty}y_n.
$$

@ID: OW23D1P3I
@BIDANG: Struktur Aljabar

Diketahui bahwa

$$
G=\{m+n\sqrt2\mid m,n\in\mathbb Z\}
$$

merupakan grup terhadap operasi penjumlahan dan

$$
H=\{3a+5b\sqrt2\mid a,b\in\mathbb Z\}
$$

merupakan subgrup dari $G$. Tentukan kardinalitas dari $G/H$.

@ID: OW23D1P4I
@BIDANG: Struktur Aljabar

Diketahui

$$
\mathbb Z_{23}\times\mathbb Z_{24}
=\{(a,b)\mid a\in\mathbb Z_{23},\,b\in\mathbb Z_{24}\}
$$

membentuk ring dengan operasi

$$
(a_1,b_1)+(a_2,b_2)
=
(a_1+a_2\!\!\!\pmod{23},\,b_1+b_2\!\!\!\pmod{24})
$$

dan

$$
(a_1,b_1)\cdot(a_2,b_2)
=
(a_1a_2\!\!\!\pmod{23},\,b_1b_2\!\!\!\pmod{24}).
$$

Tentukan banyaknya ideal di ring $\mathbb Z_{23}\times\mathbb Z_{24}$.

@ID: OW23D1P5I
@BIDANG: Kombinatorika

Diberikan 22 daerah segienam seperti pada gambar. Setiap segienam, kecuali yang diarsir hitam, akan diwarnai tepat satu warna dari tiga pilihan warna: merah, hijau, dan biru sedemikian sehingga setiap dua segienam yang memiliki sisi persekutuan berbeda warna.

Tentukan banyak cara pewarnaan yang dapat dilakukan.

@ID: OW23D1P1U
@BIDANG: Kombinatorika

Misalkan $n$ adalah bilangan bulat 8 digit dengan digit pertama adalah $5$. Berapakah peluang untuk mendapatkan $n$ yang memuat tidak lebih dari 5 digit berbeda?

@ID: OW23D1P2U
@BIDANG: Analisis Real

Diberikan barisan-barisan bilangan real $(x_n)$ dan $(y_n)$, dengan $(x_n)$ konvergen ke $x\in\mathbb R$ dan $(y_n)$ konvergen ke $y\in\mathbb R$. Jika barisan $(z_n)$ didefinisikan oleh

$$
z_n=\frac1n\sum_{i=1}^{n}x_i\,y_{n+1-i},
$$

buktikan bahwa $(z_n)$ konvergen ke $xy$.

@ID: OW23D1P3U
@BIDANG: Analisis Real

Diberikan fungsi kontinu $f:[0,1]\to\mathbb R$ yang memenuhi

$$
(f(x))^2
\le
4\int_0^x(f(t))^2\,dt
$$

untuk setiap $x\in[0,1]$.

Buktikan bahwa

$$
3(f(x))^2+2f(x)=0
$$

untuk setiap $x\in[0,1]$.

@ID: OW23D1P4U
@BIDANG: Struktur Aljabar

(a) Berikan suatu contoh homomorfisma grup $\varphi$ yang injektif dari grup $(\mathbb Z_6,+)$ ke grup $S_5$.

(b) Ada berapa banyak homomorfisma grup dari $(\mathbb Z_6,+)$ ke $S_5$? Berikan penjelasan jawaban Anda.

@ID: OW23D1P5U
@BIDANG: Struktur Aljabar

Diketahui $R$ merupakan suatu ring dengan identitas perkalian $1$ dan

$$
(x+y)^2=x^2+y^2
$$

untuk setiap $x,y\in R$.

Apakah $R$ merupakan ring komutatif? Buktikan jawaban Anda.

@ID: OW23D2P1I
@BIDANG: Kombinatorika

Tentukan banyaknya bilangan asli yang kurang dari atau sama dengan $2023$ yang memuat digit $0$.

@ID: OW23D2P2I
@BIDANG: Aljabar Linear

Misalkan $l$ adalah garis di $\mathbb R^3$ yang merupakan perpotongan bidang

$$
x+2y+3z=20
$$

dan

$$
x-y+z=23.
$$

Jika garis $l$ memotong bidang-$xy$ di titik $P(a,b,c)$, tentukan nilai $a+b+c$.

@ID: OW23D2P3I
@BIDANG: Aljabar Linear

Misalkan $I$ adalah operator identitas pada $\mathbb R^3$ dan $T:\mathbb R^3\to\mathbb R^3$ adalah operator linear yang memenuhi

$$
T\begin{pmatrix}1\\0\\0\end{pmatrix}
=
\begin{pmatrix}1\\0\\0\end{pmatrix},
\quad
T\begin{pmatrix}1\\1\\0\end{pmatrix}
=
\begin{pmatrix}1\\2\\0\end{pmatrix},
\quad
T\begin{pmatrix}1\\1\\1\end{pmatrix}
=
\begin{pmatrix}1\\2\\3\end{pmatrix}.
$$

Tentukan nilai

$$
(T-I)\circ(T-2I)\circ(T-3I)
\begin{pmatrix}9\\7\\8\end{pmatrix}.
$$

@ID: OW23D2P4I
@BIDANG: Analisis Real

Tentukan banyak bilangan kompleks tak real $z$ yang memenuhi

$$
|z-20|+|z-23|=3.
$$

@ID: OW23D2P5I
@BIDANG: Analisis Real

Tentukan nilai $m$ agar fungsi kompleks

$$
f(z)=\frac{z-2023}{1-3z}
$$

memetakan lingkaran

$$
|z-3|=m
$$

menjadi garis lurus.

@ID: OW23D2P1U
@BIDANG: Aljabar Linear

Tentukan semua bilangan real $a,b,c$ sehingga matriks

$$
A=
\begin{pmatrix}
1&2&3\\
4&5&6\\
a&b&c
\end{pmatrix}
$$

memiliki nilai eigen $2$, $0$, dan $23$.

@ID: OW23D2P2U
@BIDANG: Aljabar Linear

Misalkan $A$ matriks berukuran $n\times n$ dengan komponen real sedemikian sehingga vektor $Au$ dan $u$ ortogonal untuk setiap $u\in\mathbb R^n$.

Buktikan bahwa

$$
A^T=-A.
$$

@ID: OW23D2P3U
@BIDANG: Analisis Kompleks

Misalkan $\omega\ne1$ merupakan bilangan kompleks dengan sifat

$$
\omega^7=1.
$$

Tentukan semua bilangan bulat positif $n$ sedemikian sehingga

$$
\sum_{k=1}^{n}
\left(
1+\omega^k+\omega^{2k}+\omega^{3k}+\omega^{4k}+\omega^{5k}+\omega^{6k}
\right)
=2023.
$$

Berikan penjelasan jawaban Anda.

@ID: OW23D2P4U
@BIDANG: Analisis Real

Buktikan bahwa

$$
|\tan(x+iy)|
\ge
\frac{|e^y-e^{-y}|}{e^y+e^{-y}}.
$$

@ID: OW23D2P5U
@BIDANG: Kombinatorika

Dalam suatu turnamen terdapat $28$ tim yang akan bertanding. Masing-masing tim bermain satu sama lain tepat satu kali. Perolehan poin pada setiap pertandingan adalah $2$ poin bagi pemenang, $0$ poin bagi yang kalah, dan masing-masing $1$ poin untuk kedua tim bila berakhir seri.

Jika lebih dari $75\%$ pertandingan berakhir seri, buktikan bahwa terdapat dua tim yang meraih total poin yang sama.